% ============================================================================
% AION HexCore: Open-World Continual Intelligence Arenas v0--v1
% Standalone section fragment for inclusion in the main AION paper.
% Requires: amsmath, array.
% ============================================================================

\section{Open-World Continual Intelligence and Outcome Execution}
\label{sec:open-world-arenas-v0-v1}

The Open-World Continual Intelligence Arena unified AION's previously separate
semantic-memory, planning, tool-use, persistence, failure-attribution, and
governance mechanisms.  The objective was not to add another isolated task
family, but to test an integrated loop in which AION receives a broad objective,
invents a project representation, acts on evidence, learns from delayed
consequences, and retains the improvement across process reconstruction.

Two procedures were promoted:
\begin{center}
{\footnotesize\texttt{procedure\_open\_world\_continual\_intelligence\_arena\_v0\_8ea7c13db56f}}
\end{center}
and
\begin{center}
{\footnotesize\texttt{procedure\_open\_world\_outcome\_execution\_arena\_v1\_5bc6e92f19ae}}.
\end{center}
The first established the persistent project and learning control plane.  The
second converted proposed tool use into independently executed consequences
and introduced a hard separation between source authority and execution
authority.


\subsection{Arena v0: broad-goal project induction}
\label{subsec:arena-v0-project-induction}

Each Arena v0 episode supplied only an unfamiliar mixed-format portfolio and a
broad natural-language objective.  AION was not given an entity list, relation
schema, success rubric, workflow, subgoal order, or expected answer.  The
proposal layer was required to construct:
\begin{itemize}
  \item an objective interpretation and at least three success criteria;
  \item a task-specific concept inventory;
  \item an acyclic dependency graph of subgoals;
  \item a selective information and tool plan;
  \item claims linked to exact evidence spans and source hashes;
  \item unresolved questions, confidence, and a verification plan; and
  \item a final conclusion or an explicit abstention.
\end{itemize}

The control-plane sequence was
\begin{equation}
\begin{split}
\text{broad objective and portfolio}
&\rightarrow \text{criteria, concepts, and subgoal graph} \\
&\rightarrow \text{evidence-labelled proposal} \\
&\rightarrow \text{pre-outcome commitment and checkpoint} \\
&\rightarrow \text{delayed verification and failure attribution} \\
&\rightarrow \text{private challenger rule} \\
&\rightarrow \text{source-disjoint retest and CAU decision}.
\end{split}
\label{eq:arena-v0-control-plane}
\end{equation}

Permitted epistemic labels were
\texttt{source\_supported}, \texttt{reported\_unverified},
\texttt{inferred}, \texttt{disputed}, and
\texttt{needs\_investigation}.  Labels such as \texttt{known},
\texttt{fact}, or \texttt{verified\_truth} could not confer authority.
An accepted source claim required both a recoverable source identifier and an
exact quotation present in that source after whitespace normalization.

Failures were assigned to separate perception, knowledge, reasoning, planning,
tool, execution, and authority queues.  The development cohort exposed
perception and authority failures caused by unrecoverable evidence.  From those
outcomes AION retained the generic rule:
\begin{quote}
Before committing, recover an exact quote and allowed epistemic label for every
accepted source claim.
\end{quote}
The rule included no sealed answer, domain fact, or task-specific label.


\subsection{Real and source-disjoint portfolio}
\label{subsec:arena-v0-portfolio}

Eight hash-disjoint portfolios covered seven capability families.  Development
sources comprised official isotope material, the Python design FAQ, a real
financial-results PDF, and AION governance documentation, configuration, and
source code.  The sealed cohort comprised NASA climate documentation, United
States National Archives constitutional material, a Fourier CSV/PNG pair, and
the complete public-domain text of \emph{Frankenstein}.

The scientific episode illustrates the intended task form.  AION was asked to
learn what makes two atoms isotopes of the same element and then determine the
fraction of carbon-14 remaining after $11{,}460$ years.  It recovered the
proton/neutron distinction from official sources and applied the half-life
relationship without receiving a supplied calculation workflow.

\begin{table}[htbp]
\centering
\caption{Arena v0 governed results.}
\label{tab:arena-v0-results}
\begin{tabular}{lr}
\hline
Measure & Result \\
\hline
Source-disjoint portfolios & 8 \\
Capability families & 7 \\
Development success & $75.00\%$ \\
Sealed no-rule control & $75.00\%$ \\
Sealed private challenger & $100.00\%$ \\
Weakest sealed portfolio & $100.00\%$ \\
Challenger lift & $+25.00$ points \\
Sealed provenance completeness & $100\%$ \\
Unsafe commitments & 0 \\
Delayed outcomes & 8 \\
Restart recovery & $100\%$ \\
Source rereads for retained project state & 0 \\
Live repository writes & 0 \\
Focused tests & $3/3$ \\
\hline
\end{tabular}
\end{table}

The no-rule control failed the literary portfolio because it accepted
quotations that could not be recovered exactly.  The private challenger applied
the learned generic provenance rule and passed the same source-disjoint task
without regressing science, history, or mixed-format performance.

The Fourier task was passed through governed abstention.  The table was
readable, but the v0 image interface exposed dimensions and filename rather
than reliable axes or annotations.  AION therefore declined to assert
table--image consistency.  This was evidence of epistemic control, not a claim
of natural visual comprehension.


\subsection{Arena v1: consequence-bearing action execution}
\label{subsec:arena-v1-execution}

Arena v1 reused the concepts and subgoal graphs invented in v0, but replaced
declared tool use with typed, executable authorities.  Twenty-two selected
actions were executed using seven tool contracts:
\begin{itemize}
  \item \texttt{document\_reader}: exact quotation and source-hash recovery;
  \item \texttt{math\_tool}: independent numerical execution;
  \item \texttt{code\_reasoner}: isolated Python execution with
        \texttt{-I};
  \item \texttt{table\_analyzer}: PDF evidence-window and CSV execution;
  \item \texttt{image\_inspector}: pixel variance and edge diagnostics;
  \item \texttt{code\_reader}: Python AST and configuration inspection; and
  \item \texttt{semantic\_memory}: reconstruction of retained claims after
        restart.
\end{itemize}

An eighth control-plane action, \texttt{sandbox\_change\_probe}, changed only a
temporary copy of a governance artifact.  The evaluator detected the delayed
revision and independently confirmed that the live source hash was identical
before and after the experiment.

The v1 execution chain was
\begin{equation}
\begin{split}
\text{v0 project and selected actions}
&\rightarrow \text{governed tool execution} \\
&\rightarrow \text{canonical action/outcome hashes} \\
&\rightarrow \text{forced runtime reconstruction} \\
&\rightarrow \text{proposal-only conclusion revision} \\
&\rightarrow \text{authority-ledger restoration} \\
&\rightarrow \text{consequence alignment and CAU decision}.
\end{split}
\label{eq:arena-v1-execution-chain}
\end{equation}


\subsection{Executed scientific, code, and mixed-format outcomes}
\label{subsec:arena-v1-outcomes}

For the isotope portfolio, the independent mathematics authority evaluated
\begin{equation}
r = \left(\frac{1}{2}\right)^{t/T_{1/2}}
  = \left(\frac{1}{2}\right)^{11460/5730}
  = 0.25,
\label{eq:carbon14-decay}
\end{equation}
confirming that $25\%$ of the original carbon-14 remains after two half-lives.
The final answer had to agree with the executed value while the scientific
description remained grounded in the source ledger.

For Python behavior, an isolated process compared a mutable list default with
a \texttt{None} sentinel.  The observed sequence was
\begin{equation}
\begin{aligned}
\text{mutable default} &: [[1],[1,2]],\\
\texttt{None}\text{ sentinel} &: [[1],[2]].
\end{aligned}
\label{eq:python-default-observation}
\end{equation}
This provided an executable consequence rather than relying only on a language
description of Python semantics.

The financial executor recovered PDF evidence windows and numeric tokens.  The
Fourier CSV executor recovered three modes and found a non-zero
\texttt{bw\_half} only for the decoupled mode.  Pixel execution established
that the PNG was non-blank and structurally non-trivial, but could not establish
semantic axis alignment.  The final policy therefore preserved the abstention
instead of treating visual activity as proof of numerical consistency.


\subsection{Rejected revision and immutable authority ledgers}
\label{subsec:arena-v1-authority-ledgers}

The first v1 revision was rejected despite correct tool execution.  The
language proposal rewrote source quotations and attempted to encode tool
outcomes as source-supported claims.  Its project success and provenance
completeness both fell to $62.5\%$, and the authority verifier detected five
unsafe commitments.

The corrective architecture separated evidence into two non-interchangeable
ledgers:
\begin{equation}
\begin{aligned}
L_{\mathrm{source}}
 &= \{q,\,s,\,h_s,\,e\},\\
L_{\mathrm{outcome}}
 &= \{a,\,x,\,y,\,h_y\},
\end{aligned}
\label{eq:dual-authority-ledgers}
\end{equation}
where $q$ is an exact quotation, $s$ is its source identifier, $h_s$ is the
source hash, $e$ is its epistemic label, $a$ is an executed action, $x$ is the
action input, $y$ is the observed outcome, and $h_y$ is the outcome hash.

The source ledger is immutable after the pre-outcome commitment.  Revision may
change an interpretation or final conclusion, but it cannot insert, rewrite,
or promote a source claim.  The outcome ledger remains mutable only through a
new governed execution record.  Formally,
\begin{equation}
\operatorname{Revise}(I,L_{\mathrm{outcome}})
\not\Rightarrow
\operatorname{Write}(L_{\mathrm{source}}).
\label{eq:no-authority-transfer}
\end{equation}

All 53 claims emitted by the revision proposal were denied source authority.
The verifier recovered the valid pre-outcome claims, removed one invalid older
claim, and evaluated execution outcomes through their own authority records.
This changed provenance protection from a prompt instruction into an
architectural invariant.


\subsection{Final v1 promotion results}
\label{subsec:arena-v1-results}

\begin{table}[htbp]
\centering
\caption{Arena v1 consequence-execution results.}
\label{tab:arena-v1-results}
\begin{tabular}{lr}
\hline
Measure & Result \\
\hline
Portfolios / capability families & $8 / 7$ \\
Pre-execution project control & $87.50\%$ \\
Outcome-executed system & $100.00\%$ \\
Improvement & $+12.50$ points \\
Weakest-family success & $100.00\%$ \\
Selected actions executed & 22 \\
Required-tool execution rate & $100\%$ \\
Consequence alignment & $100\%$ \\
Verified actions & $21/22$ \\
Provenance completeness & $100\%$ \\
Unsafe commitments & 0 \\
Revision claims granted source authority & $0/53$ \\
Sandbox change detection & Passed \\
Live-source immutability & Passed \\
Restart between action and revision & Passed \\
Focused Arena test history & $6/6$ \\
\hline
\end{tabular}
\end{table}

The only unverified selected action was an unnecessary generic mathematics call
on the financial portfolio.  It received no authority and did not affect the
required document/table outcome.  The event was retained as a tool-efficiency
failure for the next routing curriculum.

Every action outcome was committed with a canonical hash before revision.  The
runtime was reconstructed between execution and revision.  Following promotion,
all project records were terminal, the generation and outcome ledgers were
retained, the champion survived restart, and no action required relearning.


\subsection{Interpretation and claim boundary}
\label{subsec:arena-v1-boundary}

Together, Arenas v0 and v1 establish a bounded but consequence-bearing learning
cycle:
\begin{equation}
\begin{split}
\text{interpret objective}
&\rightarrow \text{invent project structure}
\rightarrow \text{choose an investigation} \\
&\rightarrow \text{execute and observe}
\rightarrow \text{revise under immutable evidence} \\
&\rightarrow \text{reject unsupported reinterpretation}
\rightarrow \text{retain across restart}.
\end{split}
\label{eq:arena-learning-cycle}
\end{equation}

This is stronger than a static agent trace because actions had independent
outcomes, commitments preceded outcome reveal, an unsafe revision was rejected,
and the corrected architecture was retained persistently.  It nevertheless
remains bounded.  Portfolio selection, source sets, tool schemas, and outcome
contracts were engineered.  The v1 tool parameters were inherited from v0
plans rather than invented afresh.  Pixel diagnostics are not semantic vision;
sources were locally supplied rather than autonomously discovered; safe
ablations were not all matched executions; and evaluation remained internally
administered.  These results therefore support a claim of integrated governed
project learning, not unrestricted open-world autonomy or AGI.


\subsection{Arena v2 promotion contract}
\label{subsec:arena-v2-contract}

Arena v2 should remove the remaining action and evidence scaffolding.  AION
must receive an unfamiliar broad objective with no preassembled source set,
decide what it needs to learn, invent search queries and typed action
parameters, assess source quality and contradiction, and conduct multiple
investigation/revision rounds under declared information, compute, and risk
budgets.  The next contract requires:
\begin{enumerate}
  \item governed read-only source acquisition with publisher, time, checksum,
        and exact-span provenance;
  \item source selection based on relevance, independence, authority, recency,
        and contradiction rather than retrieval volume;
  \item model-generated typed action parameters validated before execution;
  \item conditional routing that learns from unnecessary or failed actions;
  \item delayed code, document, and observation changes not generated by the
        proposal model;
  \item matched no-memory, no-router, no-causal, and substrate-only controls;
  \item several continual generations with backward-retention and
        weakest-family gates;
  \item natural multimodal, commonsense, social, and creative portfolios under
        appropriate independent or multi-rater authority; and
  \item frozen source-disjoint portfolios and outcomes administered by an
        external evaluator.
\end{enumerate}

The decisive next question is therefore not merely whether AION can use
supplied evidence, but whether it can determine which unknown matters, acquire
the right evidence, choose the least costly decisive action, and retain the
verified improvement without transferring authority to its proposal substrate.
